\respcap{cap:gulosos}
\resp{gm-troco} Ver a §\ref{sec:guloso}: em ordem decrescente, use o máximo possível de cada moeda. Para $\{10,5,2,1\}$, uma solução ótima tem no máximo uma moeda de 1, uma de 5 e duas de 2, e nunca $2+2+1$; então as moedas pequenas somam no máximo 9, e o número de moedas de 10 é $\floor{S/10}$, como no guloso; o mesmo argumento fixa as demais. Contraexemplo geral: $\{4,3,1\}$ com $S=6$ (o guloso dá $4+1+1$; o ótimo é $3+3$).

\resp{gm-i3} \emph{Linear} e \emph{gráfica}: ver a §\ref{sec:matroides} (dimensão do espaço gerado; número de componentes de uma floresta). \emph{Transversal}: sejam $I_1,I_2\subseteq X$ emparelháveis, com $|I_1|<|I_2|$, por emparelhamentos $M_1$ e $M_2$ que cobrem, do lado $X$, exatamente $I_1$ e $I_2$. A diferença simétrica $M_1\triangle M_2$ é uma união de caminhos e ciclos alternantes; como $|M_2|>|M_1|$, algum caminho tem mais arestas de $M_2$, e começa e termina em vértices cobertos só por $M_2$. Esse caminho tem um número ímpar de arestas e alterna os lados, então um dos extremos está em $X$; ele pertence a $I_2\setminus I_1$. Trocando as arestas do caminho em $M_1$, obtém-se um emparelhamento que cobre $I_1$ e mais esse vértice. \emph{Emparelhamento}: o mesmo argumento, com caminhos alternantes num grafo qualquer: o caminho com uma aresta de $M_2$ a mais acrescenta à cobertura um vértice de $I_2$ sem descobrir os de $I_1$.

\resp{gm-g} Ver a §\ref{sec:matroides}. Suponha $|I_1|=k<|I_2|$ e que nenhum $e\in I_2\setminus I_1$ aumente $I_1$. Com pesos $k+2$ em $I_1$, $k+1$ em $I_2\setminus I_1$ e 0 nos demais, o guloso escolhe $I_1$ primeiro e depois só elementos de peso 0, totalizando $k(k+2)$; mas $I_2$ pesa pelo menos $(k+1)^2>k(k+2)$, uma contradição.

\resp{gm-pratica} (a) Não diretamente: os conjuntos de intervalos dois a dois disjuntos não formam uma matroide (com $I_1=\{[1,10]\}$ e $I_2=\{[1,2],[3,4]\}$, (I3) falha). Mesmo assim, o guloso ``o que termina primeiro'' é ótimo para pesos unitários, por um argumento de troca próprio (Exercício~\ref{ex:gm-ativ}); com pesos arbitrários, é preciso programação dinâmica (Exercício~\ref{ex:gm-pesado}). (b) Escalonamento de Gauss na matriz $m\times k$ formada pelas colunas de $S$: elas são linearmente independentes se e somente se o posto é $k$. O custo é $O(m\,k\min\{m,k\})$; se $k>m$, elas já são dependentes.

\resp{gm-pesq} A §\ref{sec:matroides} e o quadro ``Para saber mais'' reúnem os pontos centrais: os axiomas; o teorema do guloso, provado pela troca de bases; o algoritmo de Kruskal na matroide gráfica; a interseção de duas matroides por caminhos aumentantes no grafo de troca; o emparelhamento bipartido e a floresta colorida como interseções de duas matroides; e o caminho hamiltoniano como interseção de três. Para aprofundar, consulte as obras de Oxley e de Schrijver listadas nas Referências.

\resp{gm-ativ} Em ordem de término, o guloso escolhe as atividades $1,4,8,11$, isto é, os intervalos $[1,4)$, $[5,7)$, $[8,11)$ e $[12,16)$: quatro atividades, que é o máximo possível. Prova pelo argumento de troca: seja $O$ uma solução ótima e $g$ a atividade que termina primeiro. Se $g\notin O$, troque a primeira atividade de $O$ por $g$; como $g$ termina antes, continua compatível com as demais, e $|O|$ não muda. Logo existe uma solução ótima que contém $g$, e o restante é o mesmo problema restrito às atividades que começam depois do fim de $g$; por indução, o guloso é ótimo.

\resp{gm-ruins} Menor duração: com $[1,5)$, $[4,7)$ e $[6,10)$, o guloso escolhe $[4,7)$, de duração 3, e perde as outras duas, que são compatíveis entre si (o ótimo é 2). Começa primeiro: com $[0,10)$, $[1,2)$ e $[3,4)$, o guloso escolhe $[0,10)$ e fica com uma atividade; o ótimo é 2.

\resp{gm-frac} Ordene os itens por $v_i/p_i$ decrescente e leve o máximo possível de cada um. Pela troca: se uma solução ótima leva menos do item $i$ de maior razão e algo de um item $j$ de razão menor, trocar uma quantidade $\varepsilon$ de peso de $j$ por $i$ não diminui o valor; repetindo, chega-se à solução gulosa. No exemplo: os itens 1 e 2 inteiros e $20/30$ do item 3, valendo $60+100+80=$ \textbf{240}. Na mochila 0-1, o ótimo é 220 (itens 2 e 3), e o guloso pela razão dá só 160.

\resp{gm-unif} (I1) e (I2) são imediatos: as condições $|I|\le k$ e $|I\cap E_i|\le k_i$ valem para $\emptyset$ e passam para os subconjuntos. (I3) na uniforme: se $|I_1|<|I_2|\le k$, qualquer $e\in I_2\setminus I_1$ serve, pois $|I_1\cup\{e\}|\le k$. (I3) na de partição: como $|I_1|<|I_2|$, há uma parte $E_i$ com $|I_2\cap E_i|>|I_1\cap E_i|$; tome $e\in(I_2\setminus I_1)\cap E_i$. Então $|(I_1\cup\{e\})\cap E_i|\le|I_2\cap E_i|\le k_i$, e as outras partes não mudam.

\resp{gm-tarefas} Minimizar a multa das tarefas atrasadas é o mesmo que maximizar o peso das que cumprem o prazo, isto é, encontrar uma base de peso máximo da matroide. Em ordem decrescente de multa, o guloso aceita as tarefas 1, 2, 3 e 4, rejeita a 5 e a 6 (cada uma tornaria o conjunto impossível de escalonar) e aceita a 7. A multa total é $w_5+w_6=30+20=$ \textbf{50}. Uma ordem válida de execução: as tarefas aceitas em ordem crescente de prazo ($2,4,1,3,7$), seguidas das atrasadas.

\resp{gm-huffman} As junções são: $f+e=14$; $c+b=25$; $14+d=30$; $25+30=55$; $a+55=100$. Um código possível: $a=0$, $c=100$, $b=101$, $f=1100$, $e=1101$, $d=111$. O custo é $45\cdot1+13\cdot3+12\cdot3+16\cdot3+9\cdot4+5\cdot4=$ \textbf{224}.

\resp{gm-spt} Se, numa ordem, uma tarefa $i$ vem imediatamente antes de uma tarefa $j$ com $t_i>t_j$, trocar as duas não altera os términos das demais, e a soma dos términos diminui de $t_i-t_j$. Logo, numa ordem ótima, não há tal par, e as durações aparecem em ordem crescente. Para as durações $3,1,2$, a ordem $1,2,3$ termina as tarefas nos instantes $1,3,6$, com soma \textbf{10}; a pior ordem, $3,2,1$, soma 14.

\resp{gm-naomat} Não é matroide: com $I_1=\{3\}$ e $I_2=\{1,2\}$, nem $\{1,3\}$ nem $\{2,3\}$ são independentes, e (I3) falha. Com $c(3)=3$ e $c(1)=c(2)=2$, o guloso escolhe 3 e não consegue acrescentar mais nada, ficando com peso 3, enquanto $\{1,2\}$ pesa 4.

\resp{gm-paralelos} É uma matroide: é a matroide linear das colunas $1=(1,0)$, $2=(1,0)$, $3=(0,1)$ e $4=(1,1)$ de $\R^2$. Os pares independentes são todos, exceto $\{1,2\}$, cujas colunas são iguais (elementos ``paralelos''), e nenhum trio é independente em $\R^2$.

\resp{gm-posto} Uma base é uma floresta geradora maximal, isto é, uma árvore geradora em cada componente; uma componente com $n_i$ vértices contribui com $n_i-1$ arestas, e o total é $\sum_i(n_i-1)=n-c$. No exemplo, $n=6$ e $c=2$: o posto é 4 (duas arestas do triângulo e as duas do caminho).

\resp{gm-dual} As bases de $U_{n,k}$ são todos os subconjuntos de tamanho $k$; os seus complementos são todos os subconjuntos de tamanho $n-k$, que são exatamente as bases de $U_{n,n-k}$.

\resp{gm-unifgul} Na uniforme $U_{n,k}$, o guloso escolhe os $k$ elementos mais pesados. Na de partição, escolhe os $k_i$ mais pesados de cada parte $E_i$. No exemplo: $9$ e $7$ de $A$, $8$ e $3$ de $B$, $6$ e $5$ de $C$, com peso total 38.

\resp{gm-moedas} Com $S=14$, o guloso dá $10+1+1+1+1$ (5 moedas), e o ótimo é $7+7$ (2 moedas). O sistema $\{1,5,10,25\}$ não tem esse problema: o guloso é ótimo para todo valor, o que se prova com um argumento como o do exemplo da §\ref{sec:guloso}.

\resp{gm-pesado} Com $A=[0,3)$ de peso 1 e $B=[1,5)$ de peso 10, o guloso escolhe $A$, que termina primeiro, e não pode escolher $B$, ficando com peso 1; o ótimo é 10. O problema se resolve por programação dinâmica: ordene os intervalos por término, seja $p(j)$ o último intervalo compatível com o $j$-ésimo e $OPT(j)=\max\{OPT(j-1),\ w_j+OPT(p(j))\}$, em $O(n\log n)$.

\resp{gm-fracn} Encontre, por seleção em $O(n)$, a razão mediana $r$ dos itens e some os pesos dos itens com razão maior que $r$. Se essa soma é no máximo a capacidade restante, leve todos esses itens inteiros e continue com os itens de razão no máximo $r$ e a capacidade que sobra; senão, descarte os itens de razão no máximo $r$ e continue com os de razão maior. Cada passo descarta metade dos itens: $T(n)=T(n/2)+O(n)=O(n)$. No fim, sobra no máximo um item, levado fracionado.
